\subsection{Problemstellung}
Trotz der genannten vielversprechenden Potenziale von \ac{KI} in \ac{ERP}-Systemen wird der Einsatz von \ac{KI}-Agenten in der Praxis bislang nur in geringem Umfang umgesetzt. Ein wesentlicher Grund hierfür ist, dass entsprechende technische Möglichkeiten erst seit relativ kurzer Zeit verfügbar sind. Beispielsweise ist Microsoft Copilot Studio, eine Entwicklungsumgebung von Microsoft zur Erstellung von \ac{KI}-Copiloten und -Agenten, erst seit 2023 allgemein verfügbar.\cite{spataro_announcing_2023}

Darüber hinaus existieren zwar bereits wissenschaftliche Arbeiten zum Einsatz von \ac{KI} und KI-Agenten in \ac{ERP}-Systemen, diese fokussieren sich jedoch häufig auf konzeptionelle oder systemunabhängige Betrachtungen. Untersuchungen, die die Implementierung und den Einsatz von KI-Agenten in konkreten ERP-Standardlösungen, insbesondere in \ac{D365FnSCM} analysieren, sind bislang nur vereinzelt vorhanden.